\section{Conclusion}
\label{sec:conclusion}

This paper answered the challenge with one idea and the evidence to judge
it. The idea is that a project's resolved history is a localization
signal: past incidents whose reports resemble a new one vote for the
files their fixes touched, behind a firewall that admits only what an
engineer could have searched at report time. To measure it, we mined the
first public benchmark that pairs issue reports, logs, code snapshots,
and ground-truth faulty files: \nInstances{} real incidents from four
organizations, stratified by whether the report names any faulty file.

On this benchmark, fusing incident memory with log-to-code retrieval
beats the strongest baseline on all four projects on the challenge's
metrics, ranks a correct module first for 67\% to 86\% of instances
(well above random and defect-frequency baselines), and delivers its
largest gains exactly where the challenge
locates the core difficulty: reports that never name the faulty file. We
reported with equal prominence what does not work yet. On one project,
fusion still trails a single signal on that hardest stratum, and
per-regime tuning is fragile when development data is scarce.

For practice, the result reads simply. A team that keeps an issue tracker
and links fixes to issues already owns the data this method needs. The
pipeline runs on a CPU in under a second per issue, needs no training,
and every ranked file comes with citations an engineer can check: the
mentions that matched, and the past incidents whose fixes pointed there.
The benchmark, the pipeline, and a script that traces every reported
number to its result file are released with this paper.
